% TikZ 图：霍尔效应的正载流子等效模型
\begin{tikzpicture}[scale=1.05,>=Stealth,line cap=round]
  \path[use as bounding box] (-3.35,-1.95) rectangle (3.55,1.95);

  \draw[rounded corners=3pt,thick,fill=blue!4]
    (-2.40,-0.88) rectangle (2.40,0.88);
  \foreach \x in {-1.75,-0.95,-0.15,0.65,1.45} {
    \node[font=\small,blue!70!black] at (\x,0.55) {$+$};
    \node[font=\small,blue!70!black] at (\x,-0.55) {$-$};
  }

  \draw[->,very thick,red] (-2.95,0) -- (2.95,0);
  \node[red,above,font=\small] at (1.88,0) {$v_d,\ I$};
  \node[red,font=\small] at (0,-1.30)
    {正载流子等效：$v_d$ 与约定电流 $I$ 同向};

  \draw[->,orange!85!black,very thick] (-1.32,0) -- (-1.32,0.75);
  \node[orange!85!black,left,font=\small] at (-1.32,0.38) {$F_L$};
  \draw[->,purple,very thick] (0.30,0.52) -- (0.30,-0.52);
  \node[purple,right,font=\small] at (0.30,0) {$E_H$};

  \draw[<->,thick,gray] (-2.72,-0.84) -- (-2.72,0.84);
  \node[gray,left] at (-2.72,0) {$d$};
  \draw[<->,thick,blue!70!black] (2.78,-0.80) -- (2.78,0.80);
  \node[blue!70!black,right,font=\small] at (2.78,0) {$U_H$};
  \node[blue!70!black] at (3.18,0.64) {$+$};
  \node[blue!70!black] at (3.18,-0.64) {$-$};

  \foreach \x in {-1.80,-0.90,0,0.90,1.80} {
    \node[font=\large] at (\x,1.34) {$\times$};
    \node[font=\large] at (\x,-1.76) {$\times$};
  }
  \node[font=\small] at (0,1.72) {$B$ 垂直纸面向里};
\end{tikzpicture}
